\documentclass[12pt, a4paper]{article}

% --- PACOTES ---
\usepackage[utf8]{inputenc}
\usepackage[brazil]{babel}
\usepackage{abntex2cite}
\usepackage{graphicx}
\usepackage{geometry}
\usepackage{hyperref}
\usepackage{indentfirst}
\usepackage{setspace}
\usepackage{times}
\usepackage{tikz}
\usepackage{array}

% --- MARGENS / FORMATAÇÃO ---
\geometry{
 a4paper,
 left=3cm,
 right=2cm,
 top=3cm,
 bottom=2cm
}
\onehalfspacing
\setlength{\parindent}{1.25cm}

\hypersetup{
    colorlinks=true,
    linkcolor=black,
    urlcolor=blue,
    citecolor=black
}

\begin{document}

% --- CAPA ---
\begin{titlepage}
    \centering
    {\large PLANO DE PROJETO \par}
    \vspace{1.5cm}
    {\bfseries\Large ARQUITETURA EFICIENTE PARA ESTIMAÇÃO DE POSE COMPLETA (\textit{FULL-BODY}) EM AMBIENTE VEICULAR ATRAVÉS DE \textit{LIFTING} 2D-PARA-3D \par}
    \vspace{2.5cm}

    \begin{flushleft}
        \textbf{Integrante:} Davi Baechtold Campos \\
        \textbf{Orientador:} Prof. Dr. Alceu de Souza Brito Junior \\
        \textbf{Coorientador:} Prof. Dr. Alessandro Zimmer
    \end{flushleft}

    \vfill
    \begin{flushleft}
        \textbf{Assinaturas:}\\[1.4cm]
        \makebox[9cm]{\hrulefill}\\
        \small Orientador: Prof. Dr. Alceu de Souza Brito Junior \\[1.6cm]
        \makebox[9cm]{\hrulefill}\\
        \small Coorientador: Prof. Dr. Alessandro Zimmer
    \end{flushleft}

    \vfill
    {\large Curitiba, \today \par}
\end{titlepage}

\tableofcontents
\newpage
\setcounter{page}{1}

% --- RESUMO ---
\section{RESUMO}
Este projeto propõe o desenvolvimento de uma arquitetura eficiente para
estimação de pose completa (\textit{full-body}) em ambiente veicular utilizando
técnicas de \textit{2D-to-3D lifting} baseadas em \textit{keypoints}. O
\textit{lifting} refere-se ao processo de inferir coordenadas tridimensionais
(3D) a partir de pontos bidimensionais (2D) detectados em imagens, resolvendo a
ambiguidade de profundidade através de modelos de aprendizado profundo. O foco
principal será no desenvolvimento de métodos robustos para detecção e
\textit{lifting} de \textit{keypoints} corporais, com expansão posterior para
face e mãos. A pesquisa aborda a questão fundamental de como construir uma
arquitetura computacionalmente eficiente que integre corpo, face e mãos em um
sistema unificado de monitoramento veicular. O projeto utilizará
pré-treinamento com datasets públicos, incluindo o dataset do RTMPose para
\textit{2D pose keypoints}, e validação no Drive\&Act dataset para cenários
automotivos específicos.

\vspace{0.5cm}
\noindent\textbf{Palavras-chave:} Estimação de Pose \textit{Full-Body}; \textit{Keypoints}; Arquitetura Eficiente; Ambiente Veicular; \textit{2D-to-3D Lifting}.

% --- INTRODUÇÃO ---
\section{INTRODUÇÃO}
A estimação de pose completa (\textit{full-body}) em tempo real representa um
desafio crítico para sistemas de monitoramento veicular inteligente. Diferente
de aplicações gerais, o contexto automotivo demanda eficiência computacional e
robustez a oclusões específicas do ambiente de cabine.

\subsection{Motivação}
A crescente necessidade de sistemas de segurança e monitoramento em veículos
torna a estimação de pose uma área de pesquisa relevante. A capacidade de
entender a postura dos ocupantes pode melhorar a segurança e a experiência do
usuário, sendo um componente essencial para a próxima geração de Sistemas
Avançados de Assistência ao Condutor (ADAS).

\subsection{Pergunta de Pesquisa}
\textbf{Research Question:} Como construir uma arquitetura computacionalmente eficiente para estimação de pose completa (\textit{full-body}) que integre corpo, face e mãos de forma unificada em ambiente veicular restrito?

\subsection{Objetivo Geral}
Desenvolver e avaliar uma arquitetura eficiente para estimação de pose completa
utilizando \textit{keypoints}, com foco inicial em corpo e expansão progressiva
para face e mãos.

\subsection{Objetivos Específicos}
\begin{itemize}
    \item Implementar um pipeline robusto para detecção de \textit{keypoints} corporais
          2D.
    \item Desenvolver um módulo de \textit{lifting} 2D-para-3D otimizado para corpo.
    \item Expandir a arquitetura para integração de \textit{keypoints} de face e mãos.
    \item Avaliar a eficiência computacional e a precisão da arquitetura unificada.
    \item Validar o sistema em cenários automotivos realistas utilizando o Drive\&Act
          dataset.
\end{itemize}

\subsection{Estrutura do Documento}
Este documento está estruturado da seguinte forma: a Seção \ref{sec:problema}
detalha o problema a ser resolvido. A Seção \ref{sec:estado_arte} discute o
estado da arte. A Seção \ref{sec:trabalho} descreve a solução proposta. As
tecnologias, procedimentos de teste, riscos, cronograma, contribuições e
resultados esperados são apresentados nas seções subsequentes, finalizando com
as referências.

% --- DETALHAMENTO DO PROBLEMA ---
\section{DETALHAMENTO DO PROBLEMA}
\label{sec:problema}
A estimação de pose completa em ambiente veicular apresenta desafios únicos que não são comumente encontrados em outros cenários:
\begin{itemize}
    \item \textbf{Complexidade Multimodal:} Integração eficiente de corpo (17 \textit{keypoints}), face (68+ \textit{keypoints}) e mãos (21 \textit{keypoints} cada).
    \item \textbf{Oclusões Específicas:} O volante oculta mãos e torso, o painel limita a visão corporal, e a própria postura de condução restringe os movimentos e aumenta a auto-oclusão.
    \item \textbf{Eficiência Computacional:} Necessidade de processamento em tempo real com recursos de hardware limitados, típicos de sistemas embarcados.
    \item \textbf{Variabilidade Temporal:} Gestos e mudanças posturais rápidas durante a condução exigem que o modelo tenha consistência temporal.
\end{itemize}

% --- ESTADO DA ARTE ---
\section{ESTADO DA ARTE}
\label{sec:estado_arte}
Esta seção apresenta um levantamento das soluções existentes para o problema de estimação de pose, incluindo uma análise crítica das abordagens mais relevantes.

\subsection{Soluções Existentes}
Diversas abordagens têm sido propostas para a estimação de pose, incluindo
métodos baseados em redes neurais convolucionais (CNNs) e técnicas de
aprendizado profundo. A maioria dessas soluções se concentra na detecção de
\textit{keypoints} em imagens 2D, mas enfrenta desafios em ambientes dinâmicos
e com oclusões.

\subsection{Análise Crítica}
As melhores soluções atuais, como RTMPose e LiftPose3D, demonstram eficácia em
ambientes controlados, mas apresentam limitações em cenários do mundo real,
como veículos, onde a oclusão e a variabilidade de iluminação são comuns. A
proposta deste projeto visa superar essas limitações através de uma arquitetura
unificada que integra múltiplas modalidades.

% --- TRABALHO A SER DESENVOLVIDO ---
\section{TRABALHO A SER DESENVOLVIDO}
Nesta seção, uma descrição funcional da solução proposta será feita.

\subsection{Diagrama em blocos do sistema}
A Figura \ref{fig:blocos} apresenta a visão geral do sistema e a interação
entre os módulos propostos.

\begin{figure}[h!]
    \centering
    \begin{tikzpicture}[
            node distance=2cm,
            block/.style={rectangle, draw, fill=blue!20, text width=2.5cm, text centered, minimum height=1cm},
            camera/.style={rectangle, draw, fill=green!20, text width=1.8cm, text centered, minimum height=0.8cm},
            arrow/.style={thick,->,>=stealth}
        ]
        % Câmeras
        \node[camera] (cam1) {Câmera 1};
        \node[camera, below of=cam1, node distance=1cm] (cam2) {Câmera 2};

        % Módulos
        \node[block, right of=cam1, node distance=3.5cm, yshift=-0.5cm] (mod1) {Módulo 1:\\ Aquisição};
        \node[block, right of=mod1, node distance=3.5cm] (mod2) {Módulo 2:\\ Detecção 2D};
        \node[block, right of=mod2, node distance=3.5cm] (mod3) {Módulo 3:\\ Fusão e\\ Lifting 3D};
        \node[block, right of=mod3, node distance=3.5cm] (mod4) {Módulo 4:\\ Visualização};

        % Setas
        \draw[arrow] (cam1) -- (mod1);
        \draw[arrow] (cam2) -- (mod1);
        \draw[arrow] (mod1) -- (mod2);
        \draw[arrow] (mod2) -- (mod3);
        \draw[arrow] (mod3) -- (mod4);
    \end{tikzpicture}
    \caption{Diagrama em blocos do sistema proposto.}
    \label{fig:blocos}
\end{figure}

\subsection{Arquitetura Proposta}
\begin{itemize}
    \item \textbf{Backbone compartilhado:} Extrator de features eficiente para múltiplas modalidades.
    \item \textbf{Heads especializados:} Detectores específicos para corpo, face e mãos.
    \item \textbf{Fusão temporal:} Módulo de atenção para consistência inter-frame.
    \item \textbf{Lifting unificado:} Arquitetura baseada em Transformers para lifting conjunto.
\end{itemize}

\subsection{Módulos Funcionais}
Cada módulo funcional será descrito individualmente, incluindo a implementação
e as tecnologias utilizadas.

\subsubsection{Módulo 1: Aquisição}
Responsável pela captura de imagens das câmeras. Será implementado em Python
utilizando bibliotecas como OpenCV.

\subsubsection{Módulo 2: Detecção 2D}
Utilizará modelos pré-treinados para detectar keypoints corporais em imagens
2D. A implementação será feita em PyTorch, utilizando o dataset do RTMPose para
pré-treinamento.

\subsubsection{Módulo 3: Fusão e Lifting 3D}
Integrará os dados 2D e realizará o lifting para 3D. Este módulo será
desenvolvido utilizando técnicas de aprendizado profundo.

\subsubsection{Módulo 4: Visualização}
Responsável pela apresentação dos resultados. Será implementado em uma
interface gráfica utilizando Tkinter.

% --- TECNOLOGIAS E DATASETS ---
\section{TECNOLOGIAS E DATASETS}

\subsection{Tecnologias de Implementação}
\begin{itemize}
    \item \textbf{Linguagem:} Python 3.10+.
    \item \textbf{Frameworks:} PyTorch para os modelos de \textit{deep learning}.
    \item \textbf{Bibliotecas:} OpenCV para processamento de imagem, NumPy para operações numéricas.
    \item \textbf{Modelos Base:} Arquiteturas como RTMPose, EfficientNet e Transformers serão a base para os módulos.
\end{itemize}

\subsection{Datasets para Treinamento e Validação}
\begin{itemize}
    \item \textbf{Pré-treinamento 2D:} COCO, MPII, 300W (face), FreiHAND (mãos).
    \item \textbf{Treinamento do Lifting 3D:} Human3.6M, AMASS (para dados sintéticos).
    \item \textbf{Validação Automotiva:} Drive\&Act dataset.
\end{itemize}

% --- PROCEDIMENTOS DE TESTE E VALIDAÇÃO ---
\section{PROCEDIMENTOS DE TESTE E VALIDAÇÃO}
\subsection{Métricas de Avaliação}
\begin{itemize}
    \item \textbf{Precisão:} MPJPE (Mean Per Joint Position Error) e PA-MPJPE (Procrustes Aligned MPJPE) para corpo, face e mãos separadamente.
    \item \textbf{Eficiência:} FPS (Frames Per Second), latência de inferência (ms) e uso de memória (MB).
    \item \textbf{Robustez:} Avaliação da degradação da precisão sob cenários de oclusão e variações de iluminação.
\end{itemize}

\subsection{Critérios de Aceite}
\begin{itemize}
    \item \textbf{Precisão:} Atingir um MPJPE inferior a 50mm para o corpo no dataset Drive\&Act.
    \item \textbf{Eficiência:} Alcançar uma taxa de inferência de, no mínimo, 20 FPS para o sistema \textit{full-body} em uma GPU de gama média.
\end{itemize}

% --- ANÁLISE DE RISCOS ---
\section{ANÁLISE DE RISCOS}
\begin{itemize}
    \item \textbf{Complexidade da Arquitetura (Alto):} A integração de múltiplas partes do corpo é complexa. \textbf{Mitigação:} Desenvolvimento incremental (corpo -> face -> mãos), com validação em cada etapa.
    \item \textbf{Limitações Computacionais (Médio):} Atingir tempo real pode ser desafiador. \textbf{Mitigação:} Foco em arquiteturas eficientes (e.g., EfficientNet), compartilhamento de \textit{backbone} e técnicas de otimização como quantização.
    \item \textbf{Qualidade dos Dados (Médio):} As anotações do Drive\&Act podem ter limitações. \textbf{Mitigação:} Complementação com dados sintéticos gerados e calibrados para o ambiente de cabine.
\end{itemize}

% --- CRONOGRAMA DO PROJETO ---

\section{CRONOGRAMA DO PROJETO}

Descrição de todas as fases do projeto, incluindo a sequência de
desenvolvimento de cada módulo, o tempo de desenvolvimento e as fases de teste
e documentação.

\begin{table}[!ht]

    \centering

        \caption{Cronograma detalhado do projeto.}

        \label{tab:cronograma}

        \begin{tabular}{|c|p{6cm}|c|p{2.5cm}|}

            \hline

            \textbf{Fase} & \textbf{Atividades}                      & \textbf{Período} & \textbf{Entrega}              \\ \hline

            1             & Setup + dataset prep + body 2D detector  & set/2025 -- out/2025          & Pipeline corpo 2D             \\ \hline

            2             & Body lifting 3D + baseline em Drive\&Act & nov/2025 -- dez/2025          & Resultados corpo 3D           \\ \hline

            3             & Integração face: detector + lifting      & fev/2026 -- mar/2026          & Módulo face funcional         \\ \hline

            4             & Integração mãos + arquitetura unificada  & abr/2026 -- mai/2026          & Sistema fullbody              \\ \hline

            5             & Otimização + validação + TCC final       & jun/2026 -- ago/2026        & \textbf{TCC completo}         \\ \hline

        \end{tabular}

\end{table}

% --- CONCLUSÃO ---
\section{CONCLUSÃO}
Este projeto visa desenvolver uma arquitetura eficiente para estimação de pose
completa (\textit{full-body}) em ambiente veicular, abordando um problema
relevante e inovador na área de visão computacional. A estratégia de
desenvolvimento incremental, iniciando com \textit{keypoints} corporais para
uma publicação científica e expandindo para face e mãos, juntamente com o uso
de datasets específicos como o Drive\&Act, fornece uma base sólida para o
desenvolvimento de uma solução robusta que contribuirá tanto para o avanço
científico quanto para aplicações práticas em sistemas de monitoramento
veicular.

% --- REFERÊNCIAS ---
\section{REFERÊNCIAS}
\begin{thebibliography}{99}
    \bibitem{martinez2017}
    MARTINEZ, J. et al. \textbf{A simple yet effective baseline for 3D human pose estimation}. In: ICCV. 2017.

    \bibitem{iskakov2019}
    ISKAKOV, K. et al. \textbf{Learnable triangulation of human pose}. In: ICCV. 2019.

    \bibitem{driveact2019}
    MARTIN, M. et al. \textbf{Drive\&Act: A Multi-Modal Dataset for Fine-Grained Driver Behavior Recognition}. In: ICCV. 2019.

    \bibitem{rtmpose2023}
    JIANG, T. et al. \textbf{RTMPose: Real-Time Multi-Person Pose Estimation}. arXiv:2303.07399, 2023.

    \bibitem{vaswani2017}
    VASWANI, A. et al. \textbf{Attention is all you need}. In: NeurIPS. 2017.

    \bibitem{zhang2020}
    ZHANG, C. et al. \textbf{LiftPose3D: transforming 2D poses to 3D in real-time}. arXiv:2008.02945, 2020.
\end{thebibliography}

\end{document}
